\section{Proyecciones coordinadas del estado HMT: equivalencia algebraica de pantallas y realización física en dimensión 11}
\label{sec:cierre-equivalencias-moonshine-teoria-m-20260825}

Sea \(\mathcal X_{\mathrm{HMT}}^{\mathrm{enr}}\) el dominio que conserva
fase, orientación, cociente, acarreo, frontera, incidencia, gradación y
memoria. Sus componentes excepcional y de pantalla se extraen mediante las
proyecciones canónicas
\[
 \varepsilon_{\mathrm{exc}}:
 \mathcal X_{\mathrm{HMT}}^{\mathrm{enr}}
 \longrightarrow\mathscr X_{\mathrm{exc}},
 \qquad
 \varepsilon_{\mathrm{scr}}:
 \mathcal X_{\mathrm{HMT}}^{\mathrm{enr}}
 \longrightarrow\mathscr S_{\mathrm{HMT}}.
\]
La primera alimenta la construcción ya demostrada
\(\mathcal P_{\mathrm{exc}}:\mathscr X_{\mathrm{exc}}\to V^\natural\).
La segunda alimentará el morfismo algebraico
\(\mathcal R_M^{\mathrm{alg}}\) construido más abajo. La realización conjunta
es
\[
 \mathcal P_{\mathrm{joint}}^M(x)
 =\left(
   \mathcal P_{\mathrm{exc}}\varepsilon_{\mathrm{exc}}(x),
   \mathcal R_M^{\mathrm{alg}}\varepsilon_{\mathrm{scr}}(x)
  \right).
\]
Este mapa conserva el antecedente común y mantiene distintos los codominios.

\subsection{Equivalencia exacta de las pantallas reticulares y condiciones de la realización undecadimensional}

La pantalla real parte de
\[
 V_{12}\cong V_1\oplus V_3\oplus V_{3'}\oplus V_5,
 \qquad
 P_{11}=I_{12}-\frac1{12}J_{12}.
\]
El estado \(K\) produce \(u_K=P_3P_{11}K\ne0\), la recta
\(\ell_K=\mathbb Ru_K\), y
\[
 V_{11}=\ell_K\oplus V_{10},
 \qquad \dim V_{11}=11,
 \qquad \dim V_{10}=10.
\]
El paso \(12\to11\) elimina el modo uniforme y el paso \(11\to10\) elimina
la recta polarizada por \(K\).

La pantalla reticular pertenece a otra categoría. Para
\[
 L_{26}=\Lambda_{24}\oplus U\cong II_{25,1}
\]
y un vector isotrópico primitivo \(e_+\in U\),
\[
 e_+^\perp=\Lambda_{24}\oplus\mathbb Ze_+,
 \qquad
 \boxed{e_+^\perp/\mathbb Ze_+\cong\Lambda_{24}.}
\]
La ortogonalidad y el cociente eliminan dos direcciones. Este teorema no es
una repetición de \(11\to10\) ni una identificación automática con un
espacio--tiempo físico.

\subsection{Morfismo algebraico de pantalla y equivalencia cocientada}

Sea
\[
 P_{10}:=P_{11}-\frac{u_Ku_K^{\mathsf T}}{\lVert u_K\rVert^2},
 \qquad
 \mathscr G_K:=
 \{(v_{11},v_{10})\in V_{11}\times V_{10}:
             v_{10}=P_{10}v_{11}\}.
\]
La descomposición
\(e_+^\perp=\Lambda_{24}\oplus\mathbb Ze_+\) define
\[
 q_{24}:e_+^\perp\longrightarrow\Lambda_{24},
 \qquad q_{24}(\lambda+ae_+)=\lambda.
\]
Tomemos
\[
 \mathscr S_{\mathrm{HMT}}
 :=V_{12}\times\mathscr D_0\times e_+^\perp,
 \qquad
 \mathscr S_M^{\mathrm{alg}}
 :=\mathscr G_K\times\mathscr D_0\times\Lambda_{24}.
\]

\begin{definition}[Morfismo algebraico de pantalla M]
\label{def:cierre-vtm-rm-algebraico}
Para \(d=(m,w,R_0)\in\mathscr D_0\), definimos
\[
 \boxed{
 \mathcal R_M^{\mathrm{alg}}(v,d,y)
 :=\bigl((P_{11}v,P_{10}v),d,q_{24}(y)\bigr).}
\]
\end{definition}

\begin{theorem}[Equivalencia algebraica de las pantallas cocientadas]
\label{thm:cierre-vtm-rm-algebraico}
El morfismo anterior induce un isomorfismo
\[
 \boxed{
 \overline{\mathcal R}_M^{\mathrm{alg}}:
 (V_{12}/V_1)\times\mathscr D_0
 \times(e_+^\perp/\mathbb Ze_+)
 \xrightarrow{\ \cong\ }
 \mathscr S_M^{\mathrm{alg}}.}
\]
Si
\[
 \tau_{\mathrm{src}}([v],d,[y])
 =([v],\mathsf T_0d,[y])
\]
y
\[
 \tau_{\mathrm{tar}}((v_{11},v_{10}),d,\lambda)
 =((v_{11},v_{10}),\mathsf T_0d,\lambda),
\]
entonces
\[
 \overline{\mathcal R}_M^{\mathrm{alg}}\tau_{\mathrm{src}}
 =\tau_{\mathrm{tar}}\overline{\mathcal R}_M^{\mathrm{alg}}.
\]
La energía compacta \(E_{R_0}(m,w)\) se conserva bajo este
entrelazamiento y bajo radios recíprocos.
\end{theorem}

\begin{proof}
El proyector \(P_{11}\) tiene núcleo \(V_1\) e imagen \(V_{11}\). Como
\(P_{10}P_{11}=P_{10}\), el mapa
\[
 [v]\longmapsto(P_{11}v,P_{10}v)
\]
es un isomorfismo de \(V_{12}/V_1\) sobre el grafo
\(\mathscr G_K\). El mapa \(q_{24}\) es sobreyectivo y su núcleo es
\(\mathbb Ze_+\), por lo que induce el isomorfismo
\(e_+^\perp/\mathbb Ze_+\cong\Lambda_{24}\). El factor
\(\mathscr D_0\) se transporta por la identidad. El producto de estos tres
isomorfismos es \(\overline{\mathcal R}_M^{\mathrm{alg}}\).

Los proyectores y \(q_{24}\) no actúan sobre \(d\), mientras que las dos
involuciones actúan únicamente mediante \(\mathsf T_0\). El cuadrado conmuta
por sustitución directa. La invariancia energética es el
\cref{thm:cierre-vtm-tres-formulaciones-t}.
\end{proof}

Así queda construido el mapa matemático de reducción y no sólo una
correspondencia de rangos. La pantalla \(11\to10\), el cociente
\(26\to24\) y la dualidad de radio aparecen como factores distintos del
mismo morfismo.

\subsection{Reconocimiento físico undecadimensional}

La interpretación como teoría M se separa mediante un mapa de reconocimiento
\[
 \Xi_M:\mathscr S_M^{\mathrm{alg}}
 \longrightarrow\mathcal X_{11}^{\mathrm{phys}},
\]
cuyo codominio porta al menos
\[
 g_{MN},\qquad C_{MNP},\qquad\psi_M,
\]
la acción y las transformaciones de supersimetría. Una completación
operatoria incluye
\[
 \mathfrak S_M=(\mathcal H_{\mathrm{phys}},H,Q),
 \qquad Q^2=H,
\]
y una isometría
\[
 W:\mathcal H_{\mathrm{HMT}}\longrightarrow\mathcal H_{\mathrm{phys}}
\]
que, en sus dominios comunes, satisface
\[
 W\widehat H=HW,
 \qquad W\widehat Q=QW,
 \qquad \widehat Q^{\,2}=\widehat H.
\]

\begin{theorem}[Criterio de reconocimiento físico compatible]
\label{thm:cierre-vtm-reconocimiento-fisico-m}
Si \(\Xi_M\) preserva el producto interno, transporta la dirección
\(\ell_K\) a la dirección compacta, entrelaza la dualidad y satisface las
identidades operatorias anteriores, entonces el espectro físico compacto es
invariante bajo radios recíprocos y la supercarga transportada satisface
\(Q^2=H\) sobre la imagen de \(W\).
\end{theorem}

\begin{proof}
El entrelazamiento de \(\Xi_M\) con \(\tau_{\mathrm{tar}}\), unido al
\cref{thm:cierre-vtm-rm-algebraico}, transporta la involución unitaria al
sector compacto. Para la supercarga,
\[
 Q^2W=W\widehat Q^{\,2}=W\widehat H=HW.
\]
La preservación del producto interno convierte el operador transportado en
una equivalencia unitaria sobre la imagen.
\end{proof}

El isomorfismo
\(\overline{\mathcal R}_M^{\mathrm{alg}}\) es matemático y está demostrado.
El reconocimiento por \(\Xi_M\) es una interfaz física condicionada a los
campos, la acción, las restricciones de calibre, la supersimetría, los
osciladores, las amplitudes, el control de anomalías y el límite de baja
energía de la realización elegida. No se identifica el codominio físico con
\(V^\natural\): ambos reciben información coordinada del mismo antecedente
mediante funtores distintos.

No se postula una acción literal de \(\mathbb M\) sobre cifras. Esta frase es
el control de tipado colocado al final; no reemplaza la cadena positiva que
alcanza la acción profunda de \(\mathbb M\) en \(V^\natural\).
